% Integrated working draft; autogenerated from verified standalone chapter TeX on 2026-11-12.
% !TeX program = xelatex
% CellCharter DLPFC thesis working draft: chapters 2+3+4 integrated.
% The original standalone chapter files remain unchanged; the chapter_*.tex
% fragments are frozen source snapshots and must be regenerated if originals change.
% Compiles on the user's Windows computer. The two figure PDFs are real experiment
% artifacts stored locally in the results folder, not automatically backed up to GitHub.
% This is a research thesis working draft, NOT an approved school thesis template.
\documentclass[UTF8,a4paper,12pt]{ctexart}
\usepackage[margin=2.5cm]{geometry}
\usepackage{amsmath,amssymb,booktabs,tabularx,array}
\usepackage{tikz}
\usetikzlibrary{arrows.meta,positioning}
\usepackage{float}
\usepackage{graphicx}
\usepackage{pdflscape}
\usepackage{caption}
\usepackage{hyperref}
\hypersetup{hidelinks}
\pagestyle{plain}
\setlength{\parindent}{2em}
\setlength{\parskip}{0.3em}
\renewcommand{\arraystretch}{1.18}

% Only show a placeholder if the user's genuine figure PDF cannot be found.
\newcommand{\PendingFigure}[2]{%
  \begin{figure}[htbp]
    \centering
    \fbox{\begin{minipage}[c][3.0cm][c]{0.88\linewidth}
      \centering\small #2\\[0.5em]
      \textbf{待插入真实实验图件；当前仅为占位标识。}
    \end{minipage}}
    \caption{#1}
  \end{figure}
}

\begin{document}
% Section 2 is the first included chapter. Earlier thesis chapters are not authored here.
\setcounter{section}{1}
\input{D:/paper/st-clustering-research/experiments/cellcharter_dlpfc/chapter_02_background_integrated.tex}
\clearpage
\input{D:/paper/st-clustering-research/experiments/cellcharter_dlpfc/chapter_03_methods_integrated.tex}
\clearpage
\input{D:/paper/st-clustering-research/experiments/cellcharter_dlpfc/chapter_04_results_integrated.tex}
\clearpage

% One bibliography at the end; cite keys referenced in chapter 2.
\begingroup
\raggedright
\begin{thebibliography}{99}
\bibitem{lopez2018}
Lopez R, Regier J, Cole MB, et al.
Deep generative modeling for single-cell transcriptomics.
\textit{Nature Methods}, 2018, 15: 1053--1058.
doi:10.1038/s41592-018-0229-2.

\bibitem{maynard2021}
Maynard KR, Collado-Torres L, Weber LM, et al.
Transcriptome-scale spatial gene expression in the human dorsolateral prefrontal cortex.
\textit{Nature Neuroscience}, 2021, 24: 425--436.
doi:10.1038/s41593-020-00787-0.

\bibitem{varrone2024}
Varrone M, Tavernari D, Santamaria-Martínez A, et al.
CellCharter reveals spatial cell niches associated with tissue remodeling and cell plasticity.
\textit{Nature Genetics}, 2024, 56: 74--84.
doi:10.1038/s41588-023-01588-4.

\bibitem{xu2024}
Xu H, Fu H, Long Y, et al.
Unsupervised spatially embedded deep representation of spatial transcriptomics.
\textit{Genome Medicine}, 2024, 16: 12.
doi:10.1186/s13073-024-01283-x.

\bibitem{zhang2025}
Zhang W, Zhang Z, et al.
m2ST: dual multi-scale graph clustering for spatially resolved transcriptomics.
\textit{Bioinformatics}, 2025, 41(5): btaf221.
doi:10.1093/bioinformatics/btaf221.

\bibitem{sun2026}
Sun J, Biharie K, Cai P, et al.
Beyond benchmarking: an expert-guided consensus approach to spatially aware clustering.
\textit{Nature Methods}, 2026, 23: 1814--1826.
doi:10.1038/s41592-026-03194-8.
\end{thebibliography}
\endgroup
\end{document}

